% 原书第43—50页：导言及4.1.1—4.1.3。
\chapter{学习原因—结果模型}
\label{ch:learning-cause-effect}
\providecommand{\chfouraeq}{\mathrel{\stackrel{+}{=}}}
\providecommand{\chfouraleq}{\mathrel{\stackrel{+}{\leq}}}
\newcounter{chfourtransnote}
\providecommand{\chfourtnote}[1]{\stepcounter{chfourtransnote}\begingroup
\renewcommand{\thefootnote}{译注\arabic{chfourtransnote}}\footnote{#1}\addtocounter{footnote}{-1}\endgroup}

熟悉基于条件统计独立性的观测数据因果发现方法［Pearl, 2009；Spirtes et al., 2000］的读者，可能会对这样一章感到意外：它讨论的因果推断问题只有两个可观测变量，因而不存在任何非平凡的条件独立关系。本章将介绍一些假设；在这些假设下，即使只有两个可观测变量，因果推断仍然可能实现。

其中一些假设似乎过强，难以符合现实。不过，我们应当记住，即使不涉及因果问题，基于经验数据的推断也需要很强的假设。在数据维度很高而样本量很小时，尤其如此。因此，过度简化的模型随处可见，而且已经在许多学习场景中发挥了作用。

这些假设的类型十分多样，我们也确信，这里列举的假设远非全部。目前的研究仍处于探索阶段：能够使原因与结果之间的方向具有可识别性的假设，构成了一个极其广阔的空间。我们希望本章能够启发读者，进而提出其他可用于学习因果结构的假设；当然，也希望这些假设能够符合现实。

第4.1节介绍相关假设和结构可识别性的理论结果；第4.2节说明，在数据量有限时，如何利用这些结果识别结构。

\section{结构可识别性}
\subsection{为什么需要额外假设}

第3章介绍了结构因果模型，其中的结果 $E$ 通过一个函数赋值，由原因 $C$ 计算得到。读者也许会问：数据生成过程中的这种不对称性，即由 $C$ 计算 $E$ 而非反过来计算，能否仅从 $P_{C,E}$ 中看出来？也就是说，两个变量 $X,Y$ 的联合分布 $P_{X,Y}$，是否能告诉我们，它究竟由从 $X$ 到 $Y$ 的结构因果模型生成，还是由从 $Y$ 到 $X$ 的结构因果模型生成？换言之，是否可以根据联合分布识别结构？下面这个已知结果表明，如果允许使用一般的结构因果模型，答案是否定的。

\begin{proposition}[图结构的非唯一性]
对于任意两个实值变量的联合分布 $P_{X,Y}$，都存在一个结构因果模型
\[
 Y=f_Y(X,N_Y),\qquad X\indep N_Y,
\]
其中，$f_Y$ 是可测函数，$N_Y$ 是实值噪声变量。
\end{proposition}
\begin{proof}
仿照 Peters［2012，命题2.6的证明］，定义条件累积分布函数
\[
 F_{Y\mid x}(y):=P(Y\leq y\mid X=x).
\]
再定义
\[
 f_Y(x,n_Y):=F^{-1}_{Y\mid x}(n_Y),
\]
其中，$F^{-1}_{Y\mid x}(n_Y):=\inf\{x\in\RR:F_{Y\mid x}(x)\geq n_Y\}$。
最后，令 $N_Y$ 在 $[0,1]$ 上服从均匀分布，并且与 $X$ 独立。
\end{proof}

这个结果既可以用于 $X=C,Y=E$ 的情形，也可以用于 $X=E,Y=C$ 的情形。因此，每个联合分布 $P_{X,Y}$ 在两个方向上都容许结构因果模型。正因如此，人们常常认为，仅凭被动观测，无法推断两个可观测变量之间的因果方向。第7章将说明，这一论断适用于仅依据（条件）统计独立性进行因果推断的框架［Spirtes et al., 2000；Pearl, 2009］。在该框架下，因果结构 $X\to Y$ 与 $Y\to X$ 无法区分。对于仅有的两个变量，唯一可能的（条件）独立关系只能以空集为条件；除非因果影响属于非一般情形，否则 $X$ 与 $Y$ 不会因此独立。\footnote{注意，不应将这种非一般情形称为“平凡”情形，因为非平凡的反事实影响可以与 $X\indep Y$ 同时成立，见第3.4节。}
近年来，一些方法开始利用联合分布中条件独立关系之外的信息，对上述观点提出了挑战。这些方法依赖于关于概率分布与因果性之间关系的额外假设。

第4.1节的其余部分将讨论：在什么假设下，能够从联合分布恢复图结构，即结构可识别性。随后，第4.2节介绍根据有限数据集估计图的方法，即结构识别。这些统计方法不一定要以可识别性结果的证明为设计依据。严格沿着证明思路构造的方法，在数据利用效率上往往并不理想。

\subsection{假设类型概览}

\paragraph{预先限制模型类}
区分原因与结果的一种办法，是定义一个由“特别自然”的条件分布 $P_{E\mid C}$\footnote{这里用 $P_{E\mid C}$ 简记条件分布族 $(P_{E\mid C=c})_c$，并隐含地假定存在密度；换言之，$P_{E,C}$ 关于某个乘积测度绝对连续。}
和边缘分布 $P_C$ 构成的模型类。对于其中若干模型类，已有理论结果表明：边缘分布 $P_X$ 与条件分布 $P_{Y\mid X}$ 的“一般”组合所诱导的联合分布，在交换 $X$ 与 $Y$ 后，无法再用同一模型类描述。这类结论也称为可识别性结果；第4.1节后面的内容将给出相应例子。

例如，可以通过限制函数 $f_E$ 的函数类，见式 (3.2)，以及／或者限制式 (3.1)、(3.2) 中噪声分布的类别，来定义条件分布 $P_{E\mid C}$ 和边缘分布 $P_C$ 的类别。第4.1.3—4.1.6节将讨论这一做法。从机器学习的角度看，这种方法尤为自然，因为在回归、分类等标准任务中，限制函数复杂度几乎无处不在。需要注意，通过受限函数类推断因果方向时，仍然隐含地假定噪声变量相互独立，这与结构因果模型的定义一致，见定义3.1。在这个意义上，可以说这些方法利用了图2.2所示的噪声独立性；但应牢记，只有进一步限制函数类，噪声独立性才能使因果方向具有可识别性，见命题4.1。

Sun et al.［2006］、Janzing et al.［2009b］及 Comley and Dowe［2003］还给出了其他模型类的选择。例如，Sun et al.［2006］和 Janzing et al.［2009b］考虑二阶指数模型：除配分函数项外，$P_{E\mid C}$ 的对数密度是关于 $e,c$ 的二阶多项式，而 $P_C$ 的对数密度是关于 $c$ 的二阶多项式。

这部分讨论最后引出两个问题。第一，如何定义模型类，使其能够描述现实经验数据中相当一部分情况？第二，如果一个经验分布恰好只在一个方向上容许这种模型，为什么该方向就应当是因果方向？第一个问题并非因果推断所独有：只要要用函数描述观测变量之间的关系，就必须从某个“合理”的函数类中拟合函数。第二个问题则似乎触及概率与因果性关系中最深层的难题之一。对于它，我们只能提出下面这些直观而尚不精确的想法。

先给出一种直观动机，它与通常的机器学习为什么依赖受限模型类有关。正如第1章所论述的，只要从一个较小的函数类中找到了能够拟合有限数据的模型，我们就会预期，该模型也能拟合未来观测。因此，从小模型类中找到拟合数据的模型，是向未来观测泛化的关键。从形式上看，学习因果模型与通常的学习场景存在本质差别：因果模型要描述系统在干预下的行为，而不只是描述从同一分布抽取的观测。因此，不能直接照搬统计学习理论的论证，得到一套因果关系的学习理论。尽管如此，我们相信，在允许一定错误概率的意义下，从小模型类中找到一个合适模型，意味着该模型在不同背景条件下也可能成立。我们进一步认为，能够在许多不同背景条件下成立的模型，比仅能拟合单个数据集观测的模型更有可能具有因果含义，参见第7.1.6节“不同环境”的讨论。

因此，通过限制模型类来推断原因与结果，与统计学习理论的一些思想存在朦胧的联系；至于如何精确建立这种联系，仍有待今后的研究。上述关于采用小模型类中因果模型的非形式化论证，并不是说，自然界的因果关系本身确实简单。我们是否经常能够用简单函数成功拟合数据，是另一个完全不同的问题。这里要论证的仅仅是这样一种信念：如果确有一个简单函数能够拟合数据，那么它也描述了因果关系的可能性就更大。此外，第4.1.9节还将把受限模型类与原因和机制的独立性联系起来。为后面相当形式化的推导作准备，先给出一个不太现实的玩具模型；与其将它看作严肃的实例，我们更愿意将它视为一种比喻。

\paragraph{原因与机制的独立性}
第2.1节提出，$P_C$ 与 $P_{E\mid C}$ 对应自然界中两个独立的机制。因此，它们通常不包含有关对方的信息，参见原理2.1及图2.2中间的方框。自然地，一旦提出 $P_C$ 与 $P_{E\mid C}$ 在“不包含有关对方的信息”这个意义下相互独立，就必须追问：这里指的究竟是哪一种信息？我们不能明显地将这一要求形式化为某个可以用统计独立性检验来核查的条件。这是因为，在所讨论的场景中，看到的是一个固定的联合分布 $P_{C,E}$，而不是一组分布；我们无法从中检验，假定原因的分布，与给定原因时假定结果的分布，是否以相依的方式发生变化。这也正是图2.2左侧和中间方框之间的根本区别。

要用结构因果模型的语言表述原因与机制的独立性，可以假定：原因的分布，应当独立于表示因果机制的函数及噪声分布。再次强调，这本身不是关于统计独立性的陈述。它所表达的是，$f_E$ 与 $P_{N_E}$ 不包含有关 $P_C$ 的信息，反之亦然。只有当所有从 $E$ 到 $C$ 描述 $P_{C,E}$ 的结构模型都违反这一独立性时，才能利用它进行因果推断。

第4.1.7和4.1.8节将描述两个玩具场景，在其中可以对独立与相依给出明确的定义。最后，第4.1.9节将形式化 $P_C$ 与 $P_{E\mid C}$ 的独立性；这种形式化不局限于前述两个简单场景，而适用于更一般的情况。其中，相依程度用\textbf{算法互信息（algorithmic mutual information）}衡量；这一概念以 Kolmogorov 复杂度意义下的描述长度为基础。由于 Kolmogorov 复杂度不可计算，因此，它应被看作一种哲学原理，而非一种实际方法。其实际意义有两方面：一是启发新方法的开发，二是为现有方法提供依据。例如，独立性原理可以为预先限制模型类的推断方法提供理由，第4.1.9节将给出具体例子。

为了直观理解独立性与受限模型类的关系，考虑如下思想实验。从包含 $k$ 个不同边缘分布的模型类中随机选择 $P_C$，再从另一个包含 $\ell$ 个不同条件分布的模型类中选择 $P_{E\mid C}$。这样就会诱导出 $k\ell$ 个不同的联合分布 $P_{C,E}$。在一般情形下，除非模型类的定义十分特殊，这些联合分布将产生 $k\ell>k$ 个不同的边缘分布 $P_E$，以及 $k\ell>\ell$ 个不同的条件分布 $P_{C\mid E}$。因此，$P_C$ 与 $P_{E\mid C}$ 的典型组合所诱导的联合分布 $P_{E,C}$，其“反向边缘分布与条件分布”$P_E$ 和 $P_{C\mid E}$ 将不属于原来的模型类，而需要更大的模型类来描述。换句话说，无论把可能的 $P_C$ 和 $P_{E\mid C}$ 的集合取得多大，所诱导的 $P_{C\mid E}$ 和 $P_E$ 的集合还会更大。这个思想实验依赖于从有限集合中随机选择的朴素图景，因此更像一个比喻。不过，它支持这样一种看法：只要事先选定的小模型类构造得合理，因果方向上的边缘分布与条件分布就更可能在其中得到描述。

第4.1.3—4.1.6节介绍预先限制模型类的假设，第4.1.7—4.1.9节则对独立性假设进行形式化。不过，第4.1.9节的地位比较特殊，应当把它看作一种基础原理，而不是一项独立的推断方法。

\subsection{具有非高斯加性噪声的线性模型}

具有高斯噪声的线性结构方程早已得到广泛研究。相比之下，人们直到近些年才注意到［Kano and Shimizu, 2003；Shimizu et al., 2006；Hoyer et al., 2008a］，\textbf{线性非高斯无环模型（linear non-Gaussian acyclic model，LiNGAM）}为因果推断提供了新思路。特别是，根据观测数据区分“$X$ 导致 $Y$”与“$Y$ 导致 $X$”成为可能。这里的假设是，除一个加性噪声项外，结果 $E$ 是原因 $C$ 的线性函数：
\[
 E=\alpha C+N_E,\qquad N_E\indep C,
\]
其中，$\alpha\in\RR$。这也是第4.1.4节将介绍的加性噪声模型的一个特例。下面的结果表明，这个假设足以识别原因与结果。

\begin{theorem}[线性非高斯模型的可识别性]
假设 $P_{X,Y}$ 容许线性模型
\begin{equation}
 Y=\alpha X+N_Y,\qquad N_Y\indep X,
 \tag{4.1}
\end{equation}
其中，$X,N_Y,Y$ 都是连续型随机变量。那么，存在 $\beta\in\RR$ 及随机变量 $N_X$，使得
\begin{equation}
 X=\beta Y+N_X,\qquad N_X\indep Y,
 \tag{4.2}
\end{equation}
当且仅当 $N_Y$ 和 $X$ 均服从高斯分布。\chfourtnote{原书未在本定理中明确排除 $\alpha=0$。若 $\alpha=0$，则 $X\indep Y$，取 $\beta=0,N_X=X$ 即可构造反向模型，无须高斯性。因此，本结论应在非零线性作用的情形下理解。}
\end{theorem}

% 中文 TikZ 重排：坐标及曲线提取自原 PDF，未重新模拟数据。
\begin{figure}[htbp]
\centering
\definecolor{chfourfig1c0}{rgb}{1.0,1.0,1.0}
\definecolor{chfourfig1c1}{rgb}{0.13669,0.12195,0.12529}
\definecolor{chfourfig1c2}{rgb}{0.84295,0.84816,0.85371}
\definecolor{chfourfig1c3}{rgb}{0.77444,0.78143,0.78912}
\definecolor{chfourfig1c4}{rgb}{0.18155,0.18994,0.57328}
\definecolor{chfourfig1c5}{rgb}{0.93042,0.1109,0.14145}
\begin{tikzpicture}[x=0.043cm,y=-0.043cm]
\useasboundingbox (0,0) rectangle (290,144);
\path[fill=chfourfig1c0,fill opacity=1.0,draw=none] (42.4942,100.2089) --(42.4942,12.6660) --(278.9107,12.6660) --(278.9107,100.2089) --(42.4942,100.2089);
\path[fill=none,draw=chfourfig1c1,line width=0.3985pt,draw opacity=1.0] (42.4942,94.1714) --(42.4942,33.7971);
\path[fill=none,draw=chfourfig1c1,line width=0.3985pt,draw opacity=1.0] (42.4942,94.1714) --(37.1143,94.1714);
\path[fill=none,draw=chfourfig1c1,line width=0.3985pt,draw opacity=1.0] (42.4942,63.9843) --(37.1143,63.9843);
\path[fill=none,draw=chfourfig1c1,line width=0.3985pt,draw opacity=1.0] (42.4942,33.7971) --(37.1143,33.7971);
\path[fill=none,draw=chfourfig1c1,line width=0.3985pt,draw opacity=1.0] (62.1906,100.2089) --(259.2044,100.2089);
\path[fill=none,draw=chfourfig1c1,line width=0.3985pt,draw opacity=1.0] (62.1906,100.2089) --(62.1906,105.5888);
\path[fill=none,draw=chfourfig1c1,line width=0.3985pt,draw opacity=1.0] (127.8652,100.2089) --(127.8652,105.5888);
\path[fill=none,draw=chfourfig1c1,line width=0.3985pt,draw opacity=1.0] (193.5398,100.2089) --(193.5398,105.5888);
\path[fill=none,draw=chfourfig1c1,line width=0.3985pt,draw opacity=1.0] (259.2044,100.2089) --(259.2044,105.5888);
\path[fill=none,draw=chfourfig1c2,line width=0.3985pt,draw opacity=1.0,dash pattern=on .99628pt off 2.98883pt] (62.1906,100.2089) --(62.1906,12.6660);
\path[fill=none,draw=chfourfig1c2,line width=0.3985pt,draw opacity=1.0,dash pattern=on .99628pt off 2.98883pt] (95.0279,100.2089) --(95.0279,12.6660);
\path[fill=none,draw=chfourfig1c2,line width=0.3985pt,draw opacity=1.0,dash pattern=on .99628pt off 2.98883pt] (127.8652,100.2089) --(127.8652,12.6660);
\path[fill=none,draw=chfourfig1c2,line width=0.3985pt,draw opacity=1.0,dash pattern=on .99628pt off 2.98883pt] (160.7025,100.2089) --(160.7025,12.6660);
\path[fill=none,draw=chfourfig1c2,line width=0.3985pt,draw opacity=1.0,dash pattern=on .99628pt off 2.98883pt] (193.5398,100.2089) --(193.5398,12.6660);
\path[fill=none,draw=chfourfig1c2,line width=0.3985pt,draw opacity=1.0,dash pattern=on .99628pt off 2.98883pt] (226.3671,100.2089) --(226.3671,12.6660);
\path[fill=none,draw=chfourfig1c2,line width=0.3985pt,draw opacity=1.0,dash pattern=on .99628pt off 2.98883pt] (259.2044,100.2089) --(259.2044,12.6660);
\path[fill=none,draw=chfourfig1c2,line width=0.3985pt,draw opacity=1.0,dash pattern=on .99628pt off 2.98883pt] (42.4942,94.1714) --(278.9107,94.1714);
\path[fill=none,draw=chfourfig1c2,line width=0.3985pt,draw opacity=1.0,dash pattern=on .99628pt off 2.98883pt] (42.4942,79.0778) --(278.9107,79.0778);
\path[fill=none,draw=chfourfig1c2,line width=0.3985pt,draw opacity=1.0,dash pattern=on .99628pt off 2.98883pt] (42.4942,63.9843) --(278.9107,63.9843);
\path[fill=none,draw=chfourfig1c2,line width=0.3985pt,draw opacity=1.0,dash pattern=on .99628pt off 2.98883pt] (42.4942,48.8907) --(278.9107,48.8907);
\path[fill=none,draw=chfourfig1c2,line width=0.3985pt,draw opacity=1.0,dash pattern=on .99628pt off 2.98883pt] (42.4942,33.7971) --(278.9107,33.7971);
\path[fill=none,draw=chfourfig1c2,line width=0.3985pt,draw opacity=1.0,dash pattern=on .99628pt off 2.98883pt] (42.4942,18.7035) --(278.9107,18.7035);
\path[fill=chfourfig1c3,fill opacity=1.0,draw=none] (62.1906,94.1714) --(259.2044,48.8907) --(259.2044,18.7035) --(62.1906,63.9843) --(62.1906,94.1714);
\path[fill=none,draw=chfourfig1c4,line width=0.7970pt,draw opacity=1.0] (62.1906,79.0778) --(259.2044,33.7971);
\path[fill=none,draw=chfourfig1c5,line width=0.7970pt,draw opacity=1.0] (48.9402,94.1714) --(275.6629,18.7035);
\node[font=\small,inner sep=0pt,rotate=0] (ch04f1t0) at (27.0451,94.1727) {$-1$};
\node[font=\small,inner sep=0pt,rotate=0] (ch04f1t1) at (27.0451,63.9855) {$0$};
\node[font=\small,inner sep=0pt,rotate=0] (ch04f1t2) at (27.0451,33.7985) {$1$};
\node[font=\small,inner sep=0pt,rotate=0] (ch04f1t3) at (62.1903,117.0401) {$-1$};
\node[font=\small,inner sep=0pt,rotate=0] (ch04f1t4) at (127.8635,117.0401) {$0$};
\node[font=\small,inner sep=0pt,rotate=0] (ch04f1t5) at (193.5376,117.0401) {$1$};
\node[font=\small,inner sep=0pt,rotate=0] (ch04f1t6) at (259.2016,117.0401) {$2$};
\node[font=\small,inner sep=0pt,rotate=0] (ch04f1t7) at (160.2394,130.8567) {$X$};
\node[font=\small,inner sep=0pt,rotate=90] (ch04f1t8) at (12.8868,57.3129) {$Y$};
\end{tikzpicture}
\caption{一个可识别例子中 $X,Y$ 的联合密度。蓝线表示正向模型 $Y:=0.5X+N_Y$ 所对应的函数，其中 $X,N_Y$ 均服从均匀分布；灰色区域表示 $(X,Y)$ 密度的支撑集。定理4.2表明，由于 $(X,N_Y)$ 的分布为非高斯分布，不可能存在有效的反向模型。由 $(b,c)$ 刻画的红线是使 $\EE[(X-bY-c)^2]$ 最小的最小二乘拟合。它并不构成有效的反向模型 $X=bY+c+N_X$，因为所产生的噪声 $N_X$ 不会与 $Y$ 独立：$N_X$ 支撑集的大小随 $Y$ 的取值而改变。}
\label{fig:ch04-1}
\end{figure}


因此，只要 $C$ 或 $N_E$ 为非高斯变量，就足以使因果方向具有可识别性。图4.1给出了一个例子。

下面稍微深入地考察这一结果的证明。定理4.2是 Shimizu et al.［2006］提出的 LiNGAM 模型类在二元情形下的结论。该文利用\textbf{独立成分分析（independent component analysis，ICA）}［Comon, 1994，定理11］证明了定理4.2的多元版本。独立成分分析的证明以高斯分布的一种刻画为基础，这一刻画分别由 Skitovič 和 Darmois 独立证明［Skitovič, 1954, 1962；Darmois, 1953］。其结论如下。

\begin{theorem}[Darmois--Skitovič]
设 $X_1,\ldots,X_d$ 是相互独立的非退化随机变量，见附录A.1。如果存在均不为零的系数 $a_1,\ldots,a_d$ 与 $b_1,\ldots,b_d$，即对每个 $i$ 均有 $a_i\neq0\neq b_i$，使得两个线性组合
\[
 \begin{aligned}
 l_1&=a_1X_1+\cdots+a_dX_d,\\
 l_2&=b_1X_1+\cdots+b_dX_d
 \end{aligned}
\]
相互独立，那么，每个 $X_i$ 都服从正态分布。
\end{theorem}

事实上，定理4.2所述的二元版本，可以作为 Darmois--Skitovič 定理的一个直接推论，用很简短的论证证明。为便于说明，我们在附录C.1中给出了这一证明。此外，还可以证明，二元结构因果模型的可识别性能够推广到多元结构因果模型［Peters et al., 2011b］。结合这一结果，便可以从定理4.2推出 LiNGAM 的多元可识别性。

具有非高斯加性噪声的线性模型，还可以用于一个从机器学习角度看有些不同寻常、但从理论物理角度看很有意义的问题：根据数据估计时间之箭。Peters et al.［2009b］证明，自回归模型具有时间可逆性，当且仅当噪声变量服从正态分布。为探究经验时间序列的不对称性，他们拟合两个自回归模型，以此推断时间方向：一个按通常方式从过去拟合未来，另一个则从未来拟合过去。实验中，前一方向的噪声变量确实往往比逆时间方向的噪声变量更接近相互独立，参见第4.2.1节。Bauer et al.［2016］将这一思想推广到多元时间序列。Janzing［2010］把这种观测到的不对称性与热力学时间之箭联系起来。这说明，本书所讨论的原因与结果之间的不对称性，也与统计物理中的基础问题有关。
